On considère l’équation différentielle (E) : $y' + 2y = -5e^{-2x}$ où $y$ est une fonction de la variable réelle $x$, définie et dérivable sur l’ensemble $\R$ des nombres réels, et $y'$ est la fonction dérivée de $y$.

\begin{enumerate}
	\item Résoure l'équation homogène (E$_0$) : $y'+2y=0$\sautp
	
	\Lignes{2}
	\item A l’aide du logiciel Géogébra, déterminer les nombres réels $a$ et $b$ pour que la fonction $p$ définie sur $\R$ par : $p(x) = a + b~x~e^{- 2x}$ soit une solution particulière de l'équation différentielle (E), et donner l’expression de $p$. \sautm
	
	On trouve : $a=$\makebox[0.1\linewidth]{\dotfill} \hfill $b=$\makebox[0.1\linewidth]{\dotfill} \hfill $p(x)=$\makebox[0.3\linewidth]{\dotfill}
	
	\begin{tcolorbox}[colback=darkWhite,colframe=black!50!white,arc=0mm,
	left=1mm,right=1mm,top=1mm,bottom=1mm,left skip=-0.65cm]
	1) Dans le champs de saisie (en bas) : taper \verb|p(x)=a+b*x*exp(-2x)|
	
	2) mettre la courbe de $p$ en vert
	
	3) Dans le champs de saisie, taper : \verb|g(x)=p'(x)+2p(x)| {\it \footnotesize (la partie gauche de E avec $p$ au lieu de $y$)}
	
	4) mettre la courbe de $g$ en rouge
	
	5) Dans le champs de saisie, taper : \verb|d(x)= -5*exp(-2x)| {\it \footnotesize (la partie droite de E)}
	
	6) mettre $d$ en bleu
	
	7) faire bouger les curseurs (au besoin, changer les paramètres) pour faire correspondre la courbe rouge avec la courbe bleue {\it \footnotesize (donc la partie gauche avec la partie droite de E)}
	\end{tcolorbox}
	\item En déduire la forme générale des solutions de (E).\sautp
	
	\Lignes{1}
	\item Déterminer, à l'aide de géogebra et d'un curseur bien choisi, la solution particulière $f$ de (E) telle que $f(0)=1$.\sautp
	
	\Lignes{1}
	
	\begin{tcolorbox}[colback=darkWhite,colframe=black!50!white,arc=0mm,
	left=1mm,right=1mm,top=1mm,bottom=1mm,left skip=-0.65cm]
	1) Dans le champs de saisie créer une nouvelle fonction ({\bf ne pas} l'appeler $y$) comme ce que vous avez écrit à la question 3.
	
	2) la mettre en orange	
	
	3) bouger le curseur pour répondre à la question 4
	\end{tcolorbox}
\end{enumerate}